\subsection{Hilbert-Schmidt mappings}

DEFINITION 9. --- Let E and F be two hilbertian spaces. A continuous linear mapping u from E into F is called a Hilbert-Schmidt mapping if the trace of the positive endomorphism \(u^{*}u\) of E is finite. The set of all Hilbert-Schmidt mappings from E into F is denoted by \(\mathcal{L}^{2}(\mathrm{E},\mathrm{F})\) .

When \(\mathrm{E} = \mathrm{F}\) , we write \(\mathcal{L}^2(\mathrm{E})\) for \(\mathcal{L}^2(\mathrm{E};\mathrm{E})\) .

For every \(u \in \mathcal{L}(\mathrm{E},\mathrm{F})\) , let \(\| u \|_2 = \operatorname{Tr}(u^* u)^{1/2}\) , so that \(u\) belongs to \(\mathcal{L}^2(\mathrm{E};\mathrm{F})\) if and only if \(\| u \|_2\) is finite. By the definition of the trace, we get

\[
\| u \| _ {2} ^ {2} = \sup _ {x _ {1}, \dots , x _ {m}} \sum_ {i = 1} ^ {m} \left\| u (x _ {i}) \right\| ^ {2}\tag{32}
\]

where \((x_{1},\ldots ,x_{m})\) range over the set of finite orthonormal sequences in E. In particular, taking \(m = 1\) in formula (32), we have

\[
\| u \| \leqslant \| u \| _ {2} \quad (u \in \mathscr {L} (\mathrm{E}; \mathrm{F})).\tag{33}
\]

Let \((e_i)_{i\in I}\) be an orthonormal basis of \(E\) and \((f_j)_{j\in J}\) an orthonormal basis of \(F\) . By formula (25) of \(V\) , p.~50 and the Parseval's relation (V, p.~22), we have

\[
\left\| u \right\| _ {2} ^ {2} = \sum_ {i \in I} \left\| u (e _ {i}) \right\| ^ {2} = \sum_ {i, j} \left| \langle f _ {j} | u (e _ {i}) \rangle \right| ^ {2}.\tag{34}
\]

Since \(\left|\langle f_j | u(e_i) \rangle\right| = \left|\langle e_i | u^*(f_j) \rangle\right|\) , formula (34) implies that

\[
\| u ^ {*} \| _ {2} = \| u \|;\tag{35}
\]

hence, the adjoint of a Hilbert-Schmidt mapping is a Hilbert-Schmidt mapping. Let \(E_{1}\) , \(F_{1}\) be hilbertian spaces and \(v: E_{1} \to E\) , \(w: F \to F_{1}\) continuous linear mappings. From (32), we deduce immediately that

\[
\| w u \| _ {2} \leqslant \| w \|. \| u \| _ {2}.\tag{36}
\]

By (35), (36) and the relation \(uv = (v^{*}u^{*})^{*}\) , we get

\[
\left\| u v \right\| _ {2} \leqslant \left\| u \right\| _ {2} \left\| v \right\|.\tag{37}
\]

In particular, if \(u\) belongs to \(\mathcal{L}^2(\mathrm{E},\mathrm{F})\) then wuv belongs to \(\mathcal{L}^2(\mathrm{E}_1,\mathrm{F}_1)\) .

THEOREM 1. --- Let E and F be two hilbertian spaces.

\begin{enumerate}
\def\labelenumi{(\roman{enumi})}
\item
  The set \(\mathcal{L}^2 (\mathrm{E},\mathrm{F})\) is a vector subspace of \(\mathcal{L}(\mathrm{E};\mathrm{F})\) and \(u\mapsto \| u\| _2\) is a hilbertian norm (V, p.~6) on \(\mathcal{L}^2 (\mathrm{E};\mathrm{F})\) .
\item
  The isomorphism \(\theta\) from \(\mathbf{F} \otimes \mathbf{E}'\) onto \(\mathcal{L}_f(\mathbf{E};\mathbf{F})\) characterized by \(\theta(y \otimes x^*) = yx^*\) extends to an isomorphism \(\hat{\theta}\) from \(\mathbf{F} \hat{\otimes}_2 \mathbf{E}'\) onto \(\mathcal{L}^2(\mathbf{E};\mathbf{F})\) . In particular, \(\mathcal{L}_f(\mathbf{E};\mathbf{F})\) is dense in \(\mathcal{L}^2(\mathbf{E};\mathbf{F})\) .
\end{enumerate}

Let \((e_i)_{i\in \mathbf{I}}(\mathrm{resp.}(f_j)_{j\in \mathbf{J})}\) be an orthonormal basis of E (resp. F). For every \(u\in \mathcal{L}(\mathbf{E};\mathbf{F})\) let \(\Lambda (u)\) be the matrix of \(u\) with respect to chosen orthonormal bases for E and F (V, p.~22). Let \(\| \mathbf{a}\| _2\) denote the norm of an element \(a\) of the hilbertian space \(\ell^2 (\mathbf{J}\times \mathbf{I})\) . By formula (34), \(\Lambda\) is a mapping from \(\mathcal{L}^2 (\mathbf{E};\mathbf{F})\) into \(\ell^2 (\mathbf{J}\times \mathbf{I})\) such that \(\| \Lambda (u)\| _2 = \| u\|\) ; it is clear that \(\Lambda\) is injective. To prove (i), it is enough to prove that \(\Lambda\) is surjective. Let \(\mathbf{a} = (a_{ji})\) be an element of \(\ell^2 (\mathbf{J}\times \mathbf{I})\) ; by Cauchy-Schwarz inequality, we have

\[
\left| \sum_ {j, i} \overline {{{{\eta}}}} _ {j} a _ {j i} \xi_ {i} \right| ^ {2} \leqslant \sum_ {j, i} | a _ {j i} | ^ {2} \sum_ {j, i} | \overline {{{{\eta}}}} _ {j} \xi_ {i} | ^ {2} = \| \mathbf {a} \| _ {2} ^ {2} \| \boldsymbol {\xi} \| ^ {2} \| \boldsymbol {\eta} \| ^ {2}
\]

for every \(\xi = (\xi_i)\) in \(\ell^2 (\mathrm{I})\) and \(\pmb{\eta} = (\eta_j)\) in \(\ell^2 (\mathrm{J})\) . Then there exists a continuous sesquilinear form \(\Phi\) on \(\mathbf{F}\times \mathbf{E}\) such that \(\Phi (y,x) = \sum_{j,i}\overline{\eta}_ja_{ji}\xi_i\) for \(x = \sum_{i}\xi_{i}e_{i}\) in E and \(y = \sum_{j}\eta_{j}f_{j}\) in F. Let \(u\in \mathcal{L}(\mathrm{E};\mathrm{F})\) be such that \(\Phi (y,x) = \langle y|u(x)\rangle\) (V, p.~16, cor. 2). We get

\[
a _ {j i} = \Phi (f _ {j}, e _ {i}) = \langle f _ {j} | u (e _ {i}) \rangle \quad \text { for } \quad i \in \mathrm{I}, j \in \mathrm{J}  ,
\]

hence \(\mathbf{a} = \Lambda(u)\) .

Since \(\Lambda\) is a hilbertian space isomorphism from \(\mathcal{L}^2 (\mathrm{E};\mathrm{F})\) onto \(\ell^2 (\mathrm{J}\times \mathrm{I})\) and since \((f_j\otimes e_i^*)\) is an orthonormal basis of \(\mathbf{F}\hat{\otimes}_2\mathbf{E}'\) , there exists an isomorphism \(\hat{\theta}\) from \(\mathbf{F}\hat{\otimes}_2\mathbf{E}'\) onto \(\mathcal{L}^2 (\mathrm{E};\mathrm{F})\) such that

\[
\langle f _ {j} | \hat {\theta} (t) e _ {i} \rangle = \langle f _ {j} \otimes e _ {i} ^ {*} | t \rangle
\]

for every \(i \in \mathbf{I}\) , \(j \in \mathbf{I}\) and \(t \in \mathbf{F} \hat{\otimes}_2 \mathbf{E}'\) . In particular, for \(t = y \otimes x^*\) , we find

\[
\langle f _ {j} | \hat {\theta} (y \otimes x ^ {*}) e _ {i} \rangle = \langle f _ {j} \otimes e _ {i} ^ {*} | y \otimes x ^ {*} \rangle = \langle f _ {j} | y \rangle \langle x | e _ {i} \rangle = \langle f _ {j} | y x ^ {*} e _ {i} \rangle
\]

hence \(\hat{\theta}(y\otimes x^{*}) = yx^{*}\) . This proves (ii).

Q.E.D.

Examples. --- 1) Let I and J be two sets. By the proof given above, in order that a mapping \(u\) from \(\ell^2(\mathbf{I})\) into \(\ell^2(\mathbf{J})\) be a Hilbert-Schmidt mapping, it is necessary and sufficient that there exists a matrix \((a_{ji})\) in \(\ell^2(\mathbf{J} \times \mathbf{I})\) such that \(u(\xi)_j = \sum_{i \in \mathbf{I}} a_{ji} \xi_i\) for all \(\xi = (\xi_i)\) in \(\ell^2(\mathbf{I})\) .

* 2) Let X and Y be two Hausdorff topological spaces, endowed respectively with positive measures \(\mu\) and \(\nu\) . We can show that the Hilbert-Schmidt mappings from \(\mathcal{L}^2(\mathrm{X})\) into \(\mathcal{L}^2(\mathrm{Y})\) correspond bijectively to classes of square integrable functions on \(\mathrm{Y} \times \mathrm{X}\) ; to the class of a function \(\mathrm{N} \in \mathcal{L}^2(\mathrm{Y} \times \mathrm{X}, \nu \otimes \mu)\) corresponds the mapping \(u_{\mathrm{N}}\) given by

\[
(u _ {\mathrm{N}} f) (y) = \int_ {\mathrm{X}} \mathrm{N} (y, x) f (x) d \mu (x)\tag{38}
\]

for v-almost all \(y \in \mathbf{Y}\) and \(f \in \mathcal{L}^2(\mathbf{X}, \mu)\) . We have

\[
\| u _ {\mathrm{N}} \| _ {2} ^ {2} = \int_ {\mathrm{X}} \int_ {\mathrm{Y}} | \mathrm{N} (y, x) | ^ {2} d \mu (x) d v (y). *\tag{39}
\]

Remarks. --- 1) Suppose \(\mathbf{K} = \mathbf{C}\) . Let \(u\) and \(v\) be in \(\mathcal{L}^2(\mathrm{E};\mathrm{F})\) . We have the relation \(4u^{*}v = \sum_{\varepsilon^{4} = 1}\overline{\varepsilon} (u + \varepsilon v)^{*}(u + \varepsilon v)\) , hence \(u^{*}v\) belongs to \(\mathcal{L}^1(\mathrm{E})\) . The scalar product in the hilbertian space \(\mathcal{L}^2(\mathrm{E};\mathrm{F})\) is given by

\[
\langle u | v \rangle = \operatorname{Tr} (u ^ {*} v)\tag{40}
\]

since this formula defines a hermitian form on \(\mathcal{L}^{2}(\mathrm{E};\mathrm{F})\) and we get \(\langle u|u\rangle=\|u\|_{2}^{2}\) . If \(u\in\mathcal{L}^{2}(\mathrm{E};\mathrm{F})\) and \(v\in\mathcal{L}^{2}(\mathrm{F};\mathrm{E})\) , then vu belongs to \(\mathcal{L}^{1}(\mathrm{E})\) and uv to \(\mathcal{L}^{1}(\mathrm{F})\) by the preceding; moreover, we have

\[
\operatorname{Tr} (u v) = \operatorname{Tr} (v u).\tag{41}
\]

By linearity and continuity, it is enough to verify this formula when \(u = y_{1}x_{1}^{*}\) and \(v = x_{2}y_{2}^{*}\) (with \(x_{1}, x_{2}\) in E, \(y_{1}, y_{2}\) in F); but then uv is the mapping \(y \mapsto y_{1} \langle x_{1}|x_{2} \rangle \langle y_{2}|y \rangle\) and vu the mapping \(x \mapsto x_{2} \langle y_{2}|y_{1} \rangle \langle x_{1}|x \rangle\) , and (41) follows from formula (22) of V, p.~48.

Consequently, if \(u_{1}, u_{2}\) are two elements of \(\mathcal{L}^2(\mathrm{E};\mathrm{F})\) , we have, in the hilbertian space \(\mathcal{L}^2(\mathrm{F};\mathrm{E})\) ,

\[
\langle u _ {1} ^ {*} | u _ {2} ^ {*} \rangle = \operatorname{Tr} (u _ {1} u _ {2} ^ {*}) = \operatorname{Tr} (u _ {2} ^ {*} u _ {1}) = \langle u _ {2} | u _ {1} \rangle = \overline {{\langle u _ {1} | u _ {2} \rangle}};\tag{42}
\]

in other words, \(u \mapsto u^*\) is an isomorphism from the hilbertian space \(\mathcal{L}^2(\mathrm{E};\mathrm{F})\) onto the conjugate (V, p.~6) of the hilbertian space \(\mathcal{L}^2(\mathrm{F};\mathrm{E})\) . If we identify this conjugate with the dual of \(\mathcal{L}^2 (\mathrm{F};\mathrm{E})\) (V, p.~15), we see that \(\mathcal{L}^2 (\mathrm{E};\mathrm{F})\) can be identified with the dual of \(\mathcal{L}^2 (\mathrm{F};\mathrm{E})\) , the canonical bilinear form \((v,u)\mapsto \langle v,u\rangle\) being identified with \((v,u)\mapsto \operatorname {Tr}(vu)\) .

\begin{enumerate}
\def\labelenumi{\arabic{enumi})}
\setcounter{enumi}{1}
\tightlist
\item
  Suppose K = R. We leave it to the reader to verify that formulas (40) and (41) are again valid, and to show that \(\mathcal{L}^{2}(E; F)\) can be identified with the dual of \(\mathcal{L}^{2}(F; E)\) by means of the bilinear form \((u, v) \mapsto \operatorname{Tr}(uv)\) .
\end{enumerate}

